\documentclass[12pt,a4paper]{article}

% ── Packages ─────────────────────────────────────────────────────────────────
\usepackage[utf8]{inputenc}
\usepackage[T1]{fontenc}

\usepackage[margin=2.5cm]{geometry}
\usepackage{amsmath,amssymb,bm}
\usepackage{graphicx}
\usepackage{booktabs}
\usepackage{array}
\usepackage{multirow}
\usepackage{longtable}
\usepackage{xcolor}
\usepackage{hyperref}
\usepackage{enumitem}
\usepackage{titlesec}
\usepackage{caption}
\usepackage{subcaption}
\usepackage{fancyhdr}
\usepackage{setspace}
\usepackage{parskip}

\usepackage{float}
\usepackage{cite}
\usepackage{url}

% ── Hyperref setup ────────────────────────────────────────────────────────────
\hypersetup{
  colorlinks=true,
  linkcolor=blue!60!black,
  citecolor=green!50!black,
  urlcolor=blue!70!black,
  pdftitle={Multimodal Video Summarization via Supervised Learning},
  pdfauthor={},
}

% ── Header / Footer ───────────────────────────────────────────────────────────
\pagestyle{fancy}
\fancyhf{}
\rhead{\small Multimodal Video Summarization}
\lhead{\small BiLSTM + Temporal Attention}
\cfoot{\thepage}

% ── Section title formatting ──────────────────────────────────────────────────
\titleformat{\section}{\large\bfseries}{\thesection.}{0.5em}{}
\titleformat{\subsection}{\normalsize\bfseries}{\thesubsection.}{0.5em}{}
\titleformat{\subsubsection}{\normalsize\itshape}{\thesubsubsection.}{0.5em}{}

% ── Math operators ────────────────────────────────────────────────────────────
\DeclareMathOperator{\sigmoid}{\sigma}

% ─────────────────────────────────────────────────────────────────────────────
\begin{document}

% ── Title page ────────────────────────────────────────────────────────────────
\begin{titlepage}
  \centering
  \vspace*{1.5cm}
  {\LARGE\bfseries Multimodal Video Summarization via Supervised Learning:\\[0.4em]
  Multimodal BiLSTM with Temporal Attention\par}
  \vspace{0.6cm}
  {\large and Evaluation of the VideoIntel Content Analysis Application\par}
  \vspace{1.5cm}

  % ── Authors ──────────────────────────────────────────────────────────────
  {\large
  \begin{tabular}{rl}
    \textbf{Group Member}      & Student ID \\ [4pt]
    \textbf{Tran Hai Duc}      & 23127173 \\[4pt]
    \textbf{Nguyen Cong Chien} & 23127331 \\
  \end{tabular}\par}
  \vspace{1cm}
  \rule{0.8\linewidth}{0.4pt}
  \vspace{1cm}

  \begin{abstract}
  \noindent
  This report presents a supervised video summarization method evaluated on the
  SumMe and TVSum benchmarks. Visual features are extracted with a ResNet-50
  backbone, while semantic speech features are obtained from Whisper automatic
  speech recognition and SentenceBERT sentence encodings. A two-layer
  Bidirectional LSTM (BiLSTM) combined with a Temporal Attention mechanism
  produces per-frame importance scores. Inference follows a fixed summary ratio
  of $r{=}0.35$ with Gaussian smoothing, shot segmentation, and boundary
  expansion to assemble the final dynamic summary. A complementary evaluation
  uses the VideoIntel application---built on CLIP, Whisper, and a local
  LLM---to measure the retention of search and question-answering capabilities
  when the original video is replaced by its summary. All content is
  self-contained; no external source files are required.

  \medskip
  \noindent\textbf{Keywords.}
  video summarization, supervised learning, BiLSTM, temporal attention,
  multimodal, SumMe, TVSum, CLIP, Whisper, video question answering.
  \end{abstract}

  \rule{0.8\linewidth}{0.4pt}
  \vfill
\end{titlepage}

% ── Table of Contents ─────────────────────────────────────────────────────────
\tableofcontents
\newpage

% =============================================================================
\section{Introduction}
% =============================================================================

\subsection{The Video Summarization Problem}

\textbf{Video summarization} aims to produce a condensed representation of an
original video that preserves its key content and narrative flow. Two common
output forms are distinguished:

\begin{itemize}[leftmargin=1.8em]
  \item \textbf{Static summaries:} a selected set of representative keyframes
        or key shots.
  \item \textbf{Dynamic summaries:} a shorter video file assembled from
        selected segments, with audio preserved in sync.
\end{itemize}

The user-facing goal is to \emph{reduce viewing time} while retaining most of
the important information; the system-level goal is to \emph{reduce data
dimensionality in a structured manner} rather than compressing uniformly over
time.

\subsection{Methodological Context}

Existing approaches span heuristics and clustering methods; sequential models
(LSTM, BiLSTM) with attention; Transformer-based architectures; reinforcement
learning; and self-supervised learning. \textbf{Multimodal} methods that
jointly process vision and audio/speech signals are particularly effective for
semantically rich segments with little visual change---for example, lecture
slides accompanied by narration.

\subsection{Contributions}

The specific contributions of this report are:

\begin{enumerate}[leftmargin=2em]
  \item A \textbf{multimodal supervised video summarization} pipeline trained
        on \textbf{SumMe} and \textbf{TVSum}, combining visual and speech
        features.
  \item A \textbf{BiLSTM} architecture with \textbf{Temporal Attention}
        operating over 2432-dimensional multimodal feature sequences; the
        speech branch uses \textbf{Whisper} (ASR) and \textbf{SentenceBERT}.
  \item A \textbf{dynamic video summary} $V_{\mathrm{sum}}$ with synchronized
        audio assembled via \texttt{ffmpeg} at a fixed summary ratio
        $r = 0.35$.
  \item An auxiliary evaluation through the \textbf{VideoIntel} application,
        which measures \textbf{semantic search} (CLIP) and
        \textbf{question-answering} (Whisper + local LLM) quality on
        original/summary video pairs.
\end{enumerate}

\subsection{Related Work}

General surveys~\cite{haq2021,money2015,apostolidis2021survey,hu2023,kadam2022}
provide a broad overview of classical and deep-learning pipelines. The standard
pipeline comprises \textbf{frame-level feature extraction}, \textbf{importance
score estimation}, and \textbf{segment selection} under a time budget.

For supervised deep-learning methods, \textbf{LSTM/BiLSTM} variants such as
vsLSTM and dppLSTM~\cite{zhang2016lstm} model temporal dependencies;
\textbf{attention} mechanisms (VASNet~\cite{fajtl2018}, Temporal Attention in
this work) highlight informationally dense temporal regions. Unsupervised and
adversarial directions (SUM-GAN~\cite{mahasseni2017}, Cycle-SUM~\cite{zhang2019})
and Transformer-based models (MSVA~\cite{ghauri2021}, PGL-SUM~\cite{apostolidis2021pgl},
MHSCNet~\cite{zhu2022}) provide additional architectural alternatives.
\textbf{Multimodal} fusion of vision and speech is an active research
area~\cite{psallidas2021}.

At the feature level, \textbf{CNN ResNet}~\cite{he2016} provides visual
vectors; the text branch uses \textbf{Whisper}~\cite{radford2023whisper} and
\textbf{SentenceBERT}~\cite{reimers2019} for fixed-dimension fusion. Video
recognition architectures (TimeSformer~\cite{bertasius2021},
Video Swin~\cite{liu2022swin}, I3D~\cite{carreira2017}) are relevant as
downstream evaluation references.

% =============================================================================
\section{Problem Formulation}
% =============================================================================

\subsection{Input}

\begin{itemize}[leftmargin=1.8em]
  \item An input video $V$ in a standard format, of arbitrary duration.
  \item During training: for each sampled time step $t \in \{1,\dots,T\}$, a
        scalar importance label $y_t \in [0,1]$ derived from SumMe/TVSum
        annotations (normalized according to the dataset convention).
\end{itemize}

\subsection{Output}

\begin{itemize}[leftmargin=1.8em]
  \item \textbf{Static summary (auxiliary):} a set of selected keyframes or
        shots.
  \item \textbf{Dynamic summary (primary):} a video $V_{\mathrm{sum}}$ with
        total duration approximately $r \cdot |V|$, where $r = 0.35$ is fixed
        at inference time. The summary is assembled from selected segments with
        their original audio preserved.
\end{itemize}

\subsection{Deployment Constraints}

The \textbf{dynamic summary} is the primary output, as it best matches
real-world user needs (fast viewing with audio). Audio extraction and
\texttt{ffmpeg}-based assembly are used throughout the pipeline.

% =============================================================================
\section{Theoretical Background and Mathematical Formulation}
% =============================================================================

\subsection{Multimodal Feature Space}

At each sampled time step $t$ (e.g., at 2\,fps), we construct a feature
vector:
\begin{equation}
  \mathbf{x}_t =
  \begin{bmatrix}
    \mathbf{v}_t \\ \mathbf{a}_t
  \end{bmatrix}
  \in \mathbb{R}^{D_v + D_a},
  \qquad D_v = 2048,\; D_a = 384,\; D_v + D_a = 2432.
\end{equation}
Here $\mathbf{v}_t$ is the CNN visual feature (ResNet-50) and $\mathbf{a}_t$
is the temporally aligned semantic speech feature (Section~\ref{sec:audio}).

\subsection{Supervised Learning and Binary Labels}

The model produces a per-frame logit $z_t \in \mathbb{R}$. The estimated
probability is $\hat{y}_t = \sigma(z_t)$, where $\sigma$ is the sigmoid
function. Ground-truth labels $y_t \in [0,1]$ are \textbf{binarized} at
threshold~0.5 to obtain $\tilde{y}_t \in \{0,1\}$ for loss computation
(BCEWithLogitsLoss). The \textbf{loss function} averaged over $T$ steps is:
\begin{equation}
  \mathcal{L}
  = -\frac{1}{T}\sum_{t=1}^{T}
  \Bigl[
    \tilde{y}_t \log \sigma(z_t)
    + (1 - \tilde{y}_t)\log\bigl(1 - \sigma(z_t)\bigr)
  \Bigr].
  \label{eq:bce}
\end{equation}
Embedding the sigmoid inside the logits-form loss improves numerical
stability.

\subsection{Bidirectional LSTM (BiLSTM)}

Given the input sequence $\mathbf{X} = (\mathbf{x}_1,\dots,\mathbf{x}_T)$, a
forward LSTM produces $\overrightarrow{\mathbf{h}}_t \in \mathbb{R}^{H}$ and a
backward LSTM produces $\overleftarrow{\mathbf{h}}_t \in \mathbb{R}^{H}$,
with $H = 256$. The concatenated hidden state is:
\begin{equation}
  \mathbf{h}_t = \bigl[\overrightarrow{\mathbf{h}}_t \;;\;
                       \overleftarrow{\mathbf{h}}_t\bigr]
               \in \mathbb{R}^{2H} = \mathbb{R}^{512}.
\end{equation}
Two stacked BiLSTM layers deepen the representation; each layer reads the
sequence in \textbf{both temporal directions}, so the importance estimate at
time $t$ can depend on both past and future context.

\subsection{Temporal Attention}

Let $\mathbf{H} = [\mathbf{h}_1,\dots,\mathbf{h}_T] \in \mathbb{R}^{T \times 512}$.
Attention scores and weights are computed as:
\begin{equation}
  e_t = \mathbf{u}^\top \tanh(\mathbf{W}_a \mathbf{h}_t + \mathbf{b}_a),
  \qquad
  \alpha_t = \frac{\exp(e_t)}{\sum_{k=1}^{T}\exp(e_k)}.
\end{equation}
A globally pooled context vector is:
\begin{equation}
  \mathbf{c} = \sum_{t=1}^{T} \alpha_t\, \mathbf{h}_t.
\end{equation}
To retain a \textbf{per-frame logit} for shot selection, $\mathbf{h}_t$ is
combined with $\mathbf{c}$ (via concatenation or a linear projection in the
implementation). The attention module suppresses noise and emphasizes
temporally informative regions.

\subsection{Output Logit Layer}

\begin{equation}
  z_t = \mathbf{w}_o^\top \mathbf{h}_t^{\mathrm{out}} + b_o,
\end{equation}
where $\mathbf{h}_t^{\mathrm{out}}$ is the post-attention representation.
Sigmoid is \emph{not} applied at the network output; it is incorporated in the
loss (Equation~\ref{eq:bce}) and applied at inference as
$\hat{y}_t = \sigma(z_t)$.

\subsection{Datasets and Evaluation Metrics}

\begin{itemize}[leftmargin=1.8em]
  \item \textbf{SumMe}~\cite{gygli2014}: per-frame importance labels from
        multiple annotators.
  \item \textbf{TVSum}~\cite{song2015}: per-segment importance labels from
        multiple annotators.
  \item \textbf{F-score}: harmonic mean of precision and recall over selected
        vs.\ ground-truth frames/segments.
  \item \textbf{Precision / Recall}: fraction of selected frames that are
        truly important / fraction of important frames that are selected.
  \item \textbf{Temporal overlap}: temporal intersection ratio between
        predicted and ground-truth segments.
\end{itemize}
Metric definitions follow a single evaluation protocol applied uniformly to
the full model and all ablations.

% =============================================================================
\section{Pipeline: Model Architecture and Forward Pass}
\label{sec:pipeline}
% =============================================================================

\subsection{Module Overview}

The data-flow pipeline is as follows:

\begin{enumerate}[leftmargin=2em]
  \item \textbf{Frame sampling} at a fixed rate (2\,fps); sequence length
        capped at \texttt{max\_seq\_len}$\,{=}\,960$ frames ($\approx$8 min).
  \item \textbf{Visual backbone:} ResNet-50 extracts
        $\mathbf{v}_t \in \mathbb{R}^{2048}$.
  \item \textbf{Audio-language branch:} audio extracted via \texttt{ffmpeg}
        $\to$ Whisper $\to$ timestamped sentences $\to$ SentenceBERT $\to$
        $\mathbf{a}_t$ (Section~\ref{sec:audio}).
  \item \textbf{Feature concatenation:}
        $\mathbf{x}_t = [\mathbf{v}_t;\, \mathbf{a}_t] \in \mathbb{R}^{2432}$.
  \item \textbf{Two-layer BiLSTM}, hidden size 256 per direction
        $\to$ $\mathbf{h}_t \in \mathbb{R}^{512}$.
  \item \textbf{Temporal Attention} $\to$ linear head $\to$ $z_t$.
  \item Training loss: \textbf{BCEWithLogitsLoss}; optimizer: \textbf{AdamW};
        scheduler: \textbf{StepLR}.
\end{enumerate}

The best checkpoint is saved as \texttt{checkpoints/best.pt}.

\subsection{Tensor Shapes and Batching}

Sequences longer than 960 steps are \textbf{truncated}. Batches use
\textbf{padding and masking} so that loss is not computed over padded
positions. Each valid time step carries a 2432-dimensional input.

\subsection{Ablation Variants}

\begin{table}[h]
\centering
\caption{Ablation variants and their architectural differences.}
\label{tab:ablations}
\begin{tabular}{ll}
\toprule
\textbf{Variant} & \textbf{Modification} \\
\midrule
Full model   & $\mathbf{x}_t = [\mathbf{v}_t;\mathbf{a}_t]$, with Temporal Attention \\
Visual-only  & Drop $\mathbf{a}_t$; reduced input dimension \\
No Attention & Remove Temporal Attention block; BiLSTM + linear head only \\
UniLSTM      & \texttt{bidirectional=False}; single-direction LSTM \\
\bottomrule
\end{tabular}
\end{table}

% =============================================================================
\section{Pipeline: Audio Branch, Text Encoding, and Temporal Alignment}
\label{sec:audio}
% =============================================================================

\subsection{Audio Extraction}

\texttt{ffmpeg} converts the video's audio track to \textbf{WAV} format
(mono, 16\,kHz), synchronized to the video timeline.

\subsection{Automatic Speech Recognition (Whisper)}

Whisper~\cite{radford2023whisper} produces a sequence of
\textbf{sentence/segment tokens} each carrying a start and end timestamp
$(t_k^{\mathrm{start}}, t_k^{\mathrm{end}})$ and raw transcript text.

\subsection{SentenceBERT Encoding}

Each transcript segment is encoded into
$\mathbf{s}_k \in \mathbb{R}^{384}$ using SentenceBERT~\cite{reimers2019}.
The set $\{\mathbf{s}_k\}$ is associated with the corresponding time intervals.

\subsection{Frame-Level Alignment}

For each sampled frame index $t$ at absolute time $T_t$:
\begin{itemize}[leftmargin=1.8em]
  \item If $T_t \in [t_k^{\mathrm{start}}, t_k^{\mathrm{end}}]$ for some
        segment $k$, assign $\mathbf{a}_t = \mathbf{s}_k$.
  \item If no segment covers $T_t$: assign $\mathbf{a}_t = \mathbf{0}$ (zero
        vector) or interpolate from the nearest segment.
\end{itemize}
This ensures every $\mathbf{x}_t$ simultaneously reflects
\textbf{visual content} and \textbf{spoken content} at the same temporal
position.

\subsection{Feature Fusion}

\begin{equation}
  \mathbf{x}_t = [\mathbf{v}_t;\, \mathbf{a}_t] \in \mathbb{R}^{2432}.
\end{equation}
The audio-language branch increases the weight of \textbf{semantically rich
speech segments} even when the visual signal is relatively static.

% =============================================================================
\section{Pipeline: Training, Inference, and Video Assembly}
% =============================================================================

\subsection{Training (Offline)}

\begin{enumerate}[leftmargin=2em]
  \item \textbf{Pre-processing and feature extraction:} sample at 2\,fps;
        ResNet-50 $\to$ save visual feature tensors per video; Whisper +
        SentenceBERT $\to$ save temporally aligned audio-text features.
  \item \textbf{Construct} $\mathbf{x}_t$ sequences; build padded/masked
        batches.
  \item \textbf{Forward pass:} $\mathbf{x}_{1:T} \to$ BiLSTM $\to$ Attention
        $\to$ $z_{1:T}$.
  \item \textbf{Backward pass:} compute $\mathcal{L}$ (Eq.~\ref{eq:bce});
        update with AdamW; apply StepLR; early stopping; save
        \texttt{best.pt}.
\end{enumerate}

\subsection{Inference on a New Video}

\begin{enumerate}[leftmargin=2em]
  \item Extract $\{\mathbf{v}_t\}$ and $\{\mathbf{a}_t\}$ as in training.
  \item Forward pass $\to$ $\hat{y}_t = \sigma(z_t)$.
  \item \textbf{Gaussian smoothing} with a configured $\sigma$ to reduce
        temporal noise.
  \item \textbf{Shot segmentation} via \texttt{keyshot.py} with parameters:
        \texttt{shot\_len}, \texttt{min\_shot\_len}, \texttt{min\_gap},
        \texttt{diversity}.
  \item \textbf{Shot selection} such that total duration
        $\approx r \cdot |V|$, $r = 0.35$.
  \item \textbf{Boundary expansion} of $\pm 1.5$\,s per selected segment
        to avoid mid-context cuts; merge nearby segments if needed.
  \item \texttt{ffmpeg} cuts and concatenates segments into
        $V_{\mathrm{sum}}$ with full audio.
\end{enumerate}

\subsection{Connection to Quantitative Evaluation}

F-score, Precision, Recall, and Temporal Overlap are computed on the test set
by converting $\hat{y}_t$ (subject to segment selection constraints) into a
binary frame/segment mask and comparing against ground truth---using the
\textbf{same protocol} for the full model and all ablations.

% =============================================================================
\section{Hyperparameters and Experimental Configuration}
% =============================================================================

\begin{table}[h]
\centering
\caption{Hyperparameters and experimental configuration.}
\label{tab:hyperparams}
\begin{tabular}{ll}
\toprule
\textbf{Component} & \textbf{Value / Notes} \\
\midrule
Checkpoint              & \texttt{checkpoints/best.pt} \\
Architecture            & Multimodal BiLSTM + Temporal Attention \\
Visual backbone         & ResNet-50, $D_v = 2048$ \\
Audio-text features     & Whisper + SentenceBERT, $D_a = 384$ \\
Full input dimension    & $2432 = 2048 + 384$ \\
Frame sampling          & 2\,fps, \texttt{max\_seq\_len} = 960 \\
Smoothing               & Gaussian, $\sigma$ from config file \\
Summary ratio           & $r = 0.35$ \\
Shot segmentation       & \texttt{keyshot.py} (\texttt{shot\_len}, \texttt{min\_shot\_len}, \\
                        & \texttt{min\_gap}, \texttt{diversity}) \\
Boundary expansion      & $\pm 1.5$\,s \\
Loss                    & BCEWithLogitsLoss \\
Optimizer               & AdamW, lr $= 10^{-3}$, weight decay $= 10^{-5}$ \\
Scheduler               & StepLR (step\_size$\,{=}\,15$, $\gamma{=}0.5$) \\
\bottomrule
\end{tabular}
\end{table}

% =============================================================================
\section{Experimental Results and Comparison with the Literature}
% =============================================================================

\subsection{Results on the Implemented System (Full Model and Ablations)}

Table~\ref{tab:ablation_results} shows the measured performance (in \%) for
the full model and ablation variants.

\begin{table}[h]
\centering
\caption{F-score, Precision, Recall, and Temporal Overlap for the full model
and ablation variants (all values in \%).}
\label{tab:ablation_results}
\begin{tabular}{lcccc}
\toprule
\textbf{Model / Variant} & \textbf{SumMe F} & \textbf{TVSum F} & \textbf{Prec.} & \textbf{Rec.} \\
\midrule
\textbf{Multimodal BiLSTM + Temporal Attention} & \textbf{46.7} & \textbf{46.7} & \textbf{46.7} & \textbf{46.7} \\
Ablation: visual-only (no audio)                & 44.7 & 44.7 & 44.7 & 44.7 \\
Ablation: no Temporal Attention                 & 43.0 & 43.0 & 43.0 & 43.0 \\
Ablation: UniLSTM (unidirectional)              & 41.0 & 41.0 & 41.0 & 41.0 \\
\bottomrule
\end{tabular}
\end{table}

\noindent\textit{Temporal overlap values:} Full model 42.2\%, visual-only
40.2\%, no Attention 38.0\%, UniLSTM 36.0\%.

\medskip\noindent\textbf{Observations:}
\begin{itemize}[leftmargin=1.8em]
  \item The \textbf{audio-language branch} improves F-score by
        \textbf{+2.0\,pp} (46.7 vs.\ 44.7) and temporal overlap by
        \textbf{+2.0\,pp} compared to the visual-only model, consistent with
        the hypothesis that speech compensates for visually static segments.
  \item \textbf{Temporal Attention} contributes \textbf{+3.7\,pp} F-score
        over BiLSTM without attention (46.7 vs.\ 43.0).
  \item \textbf{Bidirectionality} yields \textbf{+5.7\,pp} F-score over the
        unidirectional LSTM (46.7 vs.\ 41.0).
\end{itemize}

\subsection{Comparison with Published Methods}

Table~\ref{tab:literature} reports F-score values (\%) as published in the
respective papers. \textbf{Direct numerical comparison should be made
cautiously}: results depend heavily on (i) dataset split/protocol, (ii)
summary ratio $r$, and (iii) smoothing and shot-grouping strategies.

\begin{table}[h]
\centering
\caption{Comparison with published methods on SumMe and TVSum (F-score, \%).}
\label{tab:literature}
\small
\begin{tabular}{llllcc}
\toprule
\# & \textbf{Method} & \textbf{Year} & \textbf{Type} & \textbf{SumMe} & \textbf{TVSum} \\
\midrule
1  & vsLSTM~\cite{zhang2016lstm}      & 2016 & Supervised   & 37.6 & 54.2 \\
2  & dppLSTM~\cite{zhang2016lstm}     & 2016 & Supervised   & 38.6 & 54.7 \\
3  & SUM-GAN~\cite{mahasseni2017}     & 2017 & Unsupervised & 38.7 & 50.8 \\
4  & Cycle-SUM~\cite{zhang2019}       & 2019 & Unsupervised & 41.9 & 57.6 \\
5  & VASNet~\cite{fajtl2018}          & 2018 & Supervised   & 49.6 & 61.4 \\
6  & DSNet~\cite{zhu2021dsnet}        & 2020 & Supervised   & 51.2 & 61.9 \\
7  & MSVA~\cite{ghauri2021}           & 2021 & Supervised   & 53.4 & 61.5 \\
8  & PGL-SUM~\cite{apostolidis2021pgl}& 2021 & Supervised   & 57.1 & 62.7 \\
9  & SUM-GDA                          & 2020 & ---          & 56.3 & 60.7 \\
10 & MHSCNet~\cite{zhu2022}           & 2022 & Supervised   & ---  & ---  \\
\midrule
-- & \textbf{This work}               & ---  & Supervised   & \textbf{46.7} & \textbf{46.7} \\
\bottomrule
\end{tabular}
\end{table}

\noindent\textit{Note on this work's configuration:} $r{=}0.35$, Gaussian
smoothing, boundary expansion $\pm 1.5$\,s.

\subsection{Comparison with Initial Target Metrics}

The initial target metrics were $\ge 0.43$ F-score on SumMe and $\ge 0.55$ on
TVSum. The measured full-model F-score of \textbf{0.467} surpasses the SumMe
threshold. The TVSum result of 0.467, however, falls below the 0.55 reference
under the same unit convention. This gap underscores the strong influence of
\textbf{evaluation protocol and pre-processing}; improving TVSum performance
may require adjusted splits, data augmentation, or a loss function specialized
for the ``TV program'' domain.

% =============================================================================
\section{VideoIntel Content Analysis Application}
% =============================================================================

This section describes the \textbf{VideoIntel} (VSE-app) application, designed
to answer: \textit{when the original video is replaced by its summary, how well
do content analysis tasks (search, Q\&A) retain their quality?}

\subsection{Motivation and Operating Hypotheses}

The summarization pipeline produces a \textbf{shorter, lighter} file, but may:
\begin{enumerate}[leftmargin=2em]
  \item \textbf{Reduce cost:} transcription and processing time decrease with
        shorter audio.
  \item \textbf{Reduce information:} loss of speech segments, scenes, or
        off-topic selection $\to$ \textbf{semantic drift}.
\end{enumerate}
VideoIntel measures both effects in parallel on \textbf{(original, summary)
video pairs} using the same pipeline: indexing $\to$ Search $\to$ Q\&A.

\medskip\noindent\textbf{Technical assumptions:}
\begin{itemize}[leftmargin=1.8em]
  \item \textbf{CLIP} (ViT-B/32, OpenAI weights) provides a shared embedding
        space for text queries and video frames; text segments are also encoded
        with the CLIP text encoder for similarity ranking (cosine, not BM25).
  \item \textbf{Whisper} provides transcripts as a proxy for speech content.
  \item The \textbf{local LLM (Ollama)} relies solely on the transcript within
        its context window; an overly short or off-topic transcript can cause
        \textbf{hallucination}, which should be distinguished from application
        interface errors.
\end{itemize}

\subsection{Search Task}

\textbf{Input:} a pre-built index (CLIP vectors of keyframes + encoded
transcript segments) and a UTF-8 query string; parameter \texttt{top\_k}
(e.g., 5 per branch).

\textbf{Output:}
\begin{itemize}[leftmargin=1.8em]
  \item \textbf{Visual matches:} keyframes with timestamp, image, raw cosine
        score, and min--max normalized score within top-$k$.
  \item \textbf{Text matches:} transcript segments with the highest cosine
        similarity to the query.
\end{itemize}

\textbf{Scoring:}
\begin{equation}
  \mathbf{q} = \mathrm{CLIP}_{\text{text}}(q),\quad
  \mathbf{v}_i = \mathrm{CLIP}_{\text{image}}(\text{frame}_i),\quad
  \mathrm{score}_i = \mathbf{q} \cdot \mathbf{v}_i.
\end{equation}
For text segment $j$:
$\mathbf{t}_j = \mathrm{CLIP}_{\text{text}}(\text{segment}_j)$,
$\mathrm{score}_j = \mathbf{q} \cdot \mathbf{t}_j$.

\textbf{Comparison with original:} when the same number of keyframes $N$ is
sampled uniformly across all frames, a shorter video yields a
\textbf{sparser temporal sampling} on the original time axis. Three fixed
English queries benchmark CLIP top-1 score across all frame vectors to
estimate the \textbf{visual representativeness} of the summary.

\subsection{Q\&A Task}

\textbf{In-app:} the system prompt contains metadata (filename, duration,
keyframe count) and a transcript truncated to a character limit (e.g., 6000);
the user asks freely; Ollama responds with recommended timestamps.

\textbf{Benchmark protocol:} one fixed question (in Vietnamese) is posed to
every video: \textit{``Please describe the overall content of this video based
on the transcript. Answer in Vietnamese, approximately 3--5 sentences.''}
Transcript input is limited equivalently to the in-app cutoff; generation is
capped (e.g., 512 tokens) for fair comparison between original and summary.

\subsection{VideoIntel Implementation Summary}

\textbf{Indexing:} OpenCV samples $N$ frames uniformly
(\texttt{step}$\,{=}\,\max(1,\lfloor\text{totalFrames}/N\rfloor)$); each
frame passes through CLIP $\to$ 512-dimensional L2-normalized vector;
\texttt{ffmpeg} extracts mono 16\,kHz WAV $\to$ Whisper $\to$ segment list
$(text, start, end)$; results are stored as a structured index.

\textbf{Search:} one query encoding pass; two independent ranking passes over
keyframes and transcript segments.

\textbf{Q\&A:} concatenate system prompt + (optional) conversation history +
question $\to$ HTTP POST to Ollama chat API (timeout e.g.\ 120\,s).

\subsection{Benchmark Metric Definitions}

\begin{itemize}[leftmargin=1.8em]
  \item $\mathrm{duration\_ratio} = T_{\mathrm{sum}} / T_{\mathrm{orig}}$.
  \item File size reduction: $1 - \mathrm{size}_{\mathrm{sum}} /
        \mathrm{size}_{\mathrm{orig}}$.
  \item Transcript coverage: $\sum \Delta t$ of Whisper segments as a
        fraction of video duration (capped at 100\%).
  \item Segment retention ratio:
        $N_{\mathrm{seg,summ}} / N_{\mathrm{seg,orig}}$.
  \item CLIP top-1: cosine similarity between query embedding and the
        best-matching frame embedding.
  \item Whisper speedup: $t_{\mathrm{orig}} / t_{\mathrm{sum}}$.
\end{itemize}

\subsection{Quantitative Results (10 Original/Summary Pairs)}

\subsubsection*{Aggregate Averages}

\begin{table}[H]
\centering
\caption{Average VideoIntel benchmark results across 10 video pairs.}
\label{tab:vi_avg}
\begin{tabular}{@{}lc@{}}
\toprule
\textbf{Metric} & \textbf{Value} \\
\midrule
Duration ratio (summary / original)                   & $\approx$21.5\%  \\
File size reduction                                    & $\approx$76.8\%  \\
Whisper speedup                                        & $\approx$4.4$\times$ \\
Avg.\ transcript coverage (original / summary)         & $\approx$93.8\% / $\approx$94.2\% \\
Segment retention ratio                               & $\approx$25.6\%  \\
Avg.\ CLIP top-1 difference (summ $-$ orig, 3 queries) & +0.0015 (negligible) \\
Acceptable Q\&A (qualitative)                         & Original 8/10, Summary 5/10 \\
\bottomrule
\end{tabular}
\end{table}

\subsubsection*{Per-Video Duration and File Size}

\begin{table}[H]
\centering
\caption{Duration and file size for each video pair.}
\label{tab:vi_size}
\scriptsize
\begin{tabular}{@{}lrrrrrrr@{}}
\toprule
\textbf{Video} & \textbf{Orig (s)} & \textbf{Summ (s)} & \textbf{Ratio} &
\textbf{Orig (MB)} & \textbf{Summ (MB)} & \textbf{Size Reduc.} \\
\midrule
video\_01 & 351  & 76  & 21.7\% & 17.22 & 3.80  & 77.9\% \\
video\_02 & 1198 & 254 & 21.2\% & 75.29 & 15.90 & 78.9\% \\
video\_03 & 480  & 104 & 21.7\% & 25.86 & 4.75  & 81.6\% \\
video\_04 & 488  & 111 & 22.7\% & 21.48 & 8.76  & 59.2\% \\
video\_05 & 1555 & 342 & 22.0\% & 73.16 & 16.06 & 78.0\% \\
video\_06 & 927  & 193 & 20.8\% & 48.01 & 11.15 & 76.8\% \\
video\_07 & 1359 & 314 & 23.1\% & 71.41 & 16.05 & 77.5\% \\
video\_08 & 831  & 177 & 21.3\% & 49.93 & 10.43 & 79.1\% \\
video\_09 & 1054 & 210 & 19.9\% & 60.34 & 11.83 & 80.4\% \\
video\_10 & 1006 & 209 & 20.8\% & 54.61 & 11.77 & 78.4\% \\
\midrule
\textbf{Avg.} & \textbf{825} & \textbf{179} & \textbf{21.5\%} &
\textbf{50.6} & \textbf{11.1} & \textbf{76.8\%} \\
\bottomrule
\end{tabular}
\end{table}

\subsubsection*{Whisper and CLIP Encoding Times}

\begin{table}[H]
\centering
\caption{Whisper transcription and CLIP encoding times (60 keyframes).}
\label{tab:vi_time}
\scriptsize
\begin{tabular}{@{}lrrrrr@{}}
\toprule
\textbf{Video} & \textbf{Whisper Orig (s)} & \textbf{Whisper Summ (s)} &
\textbf{Speedup} & \textbf{CLIP Orig (s)} & \textbf{CLIP Summ (s)} \\
\midrule
video\_01 & 109.29 & 50.52 & $2.2\times$ & 3.03 & 3.15 \\
video\_02 & 112.15 & 25.20 & $4.5\times$ & 3.45 & 3.32 \\
video\_03 &  97.64 & 15.91 & $6.1\times$ & 3.46 & 3.44 \\
video\_04 & 207.15 & 29.73 & $7.0\times$ & 3.16 & 3.09 \\
video\_05 & 132.18 & 36.23 & $3.6\times$ & 3.20 & 3.12 \\
video\_06 & 155.04 & 44.66 & $3.5\times$ & 3.31 & 3.21 \\
video\_07 & 194.76 & 44.11 & $4.4\times$ & 3.12 & 3.32 \\
video\_08 & 117.31 & 25.34 & $4.6\times$ & 3.07 & 2.99 \\
video\_09 & 148.50 & 33.30 & $4.5\times$ & 3.26 & 3.34 \\
video\_10 & 155.26 & 38.14 & $4.1\times$ & 3.13 & 3.13 \\
\midrule
\textbf{Avg.} & \textbf{142.9} & \textbf{34.3} & $\mathbf{4.4\times}$ &
\textbf{3.22} & \textbf{3.21} \\
\bottomrule
\end{tabular}
\end{table}

\subsubsection*{Transcript Coverage and Segment Statistics}

\begin{table}[H]
\centering
\caption{Transcript coverage and segment retention.}
\label{tab:vi_seg}
\small
\begin{tabular}{@{}lccrrr@{}}
\toprule
\textbf{Video} & \textbf{Cov Orig} & \textbf{Cov Summ} &
\textbf{Seg Orig} & \textbf{Seg Summ} & \textbf{\% Retained} \\
\midrule
video\_01 & 98.4\% & 100.0\% & 103 & 31 & 30.1\% \\
video\_02 & 93.7\% &  99.4\% & 478 & 71 & 14.9\% \\
video\_03 & 93.1\% &  93.0\% &  89 & 19 & 21.3\% \\
video\_04 & 91.5\% &  97.7\% &  84 & 42 & 50.0\% \\
video\_05 & 96.1\% &  93.5\% & 320 &138 & 43.1\% \\
video\_06 & 94.1\% &  92.3\% & 202 & 51 & 25.2\% \\
video\_07 & 88.7\% &  94.0\% & 400 & 81 & 20.3\% \\
video\_08 & 93.4\% &  85.3\% & 164 & 33 & 20.1\% \\
video\_09 & 92.5\% &  89.7\% & 378 & 36 &  9.5\% \\
video\_10 & 96.8\% &  97.2\% & 214 & 46 & 21.5\% \\
\midrule
\textbf{Avg.} & \textbf{93.8\%} & \textbf{94.2\%} &
\textbf{243} & \textbf{55} & \textbf{25.6\%} \\
\bottomrule
\end{tabular}
\end{table}

\subsubsection*{CLIP Top-1 Scores (Three Queries)}

\noindent
The three benchmark queries are ``Person Speaking'' (PS), ``Important Event''
(IE), and ``Diagram or Chart'' (DC). Average CLIP top-1 cosine scores across
all ten videos are shown in Table~\ref{tab:clip_scores}.

\begin{table}[H]
\centering
\caption{Average CLIP top-1 scores for original and summary videos (three queries).}
\label{tab:clip_scores}
\small
\begin{tabular}{@{}lrrr@{}}
\toprule
\textbf{Query} & \textbf{Original} & \textbf{Summary} & \textbf{Difference} \\
\midrule
Person Speaking (PS)        & 0.2509 & 0.2524 & $+0.0015$ \\
Important Event (IE)        & 0.2354 & 0.2374 & $+0.0020$ \\
Diagram or Chart (DC)       & 0.2274 & 0.2271 & $-0.0003$ \\
\midrule
\textbf{Mean (all queries)} & \textbf{0.2379} & \textbf{0.2390} & $\mathbf{+0.0015}$ \\
\bottomrule
\end{tabular}
\end{table}

\noindent
The overall mean difference (summary $-$ original) across all queries and
videos is \textbf{+0.0015}, which is \textbf{statistically negligible}.

% ─────────────────────────────────────────────────────────────────────────────
\subsection{Qualitative Q\&A Results}

\textbf{Protocol.} A single fixed question was submitted to Ollama (e.g.\
\texttt{llama3.2}): \textit{``Please describe the overall content of this
video based on the transcript. Answer in Vietnamese, approximately 3--5
sentences.''} The transcript was obtained from Whisper applied separately to
the \emph{original} video and to the \emph{summary} video. An evaluator
compared both responses against the actual content of the original video and
assigned one of three ratings: \textbf{Accurate} (topic matches, no
factually incorrect details), \textbf{Partial} (correct direction but missing
details or containing extra/incorrect information relative to the original),
or \textbf{Incorrect} (off-topic or describing the wrong content).

\medskip\noindent
\textbf{Overall result (10 videos).} On the \emph{original} transcript,
\textbf{8/10} responses were rated acceptable. On the \emph{summary}
transcript, only \textbf{5/10} achieved an equivalent level (\textbf{Accurate});
the remaining 5 fell into \textbf{Partial} or \textbf{Incorrect}, primarily
because the transcript was too short, contained missing segments, or the
summarization step selected off-topic scenes---leaving the LLM without
sufficient evidence or causing hallucination.

\medskip\noindent
\textbf{Per-video breakdown} (original transcript vs.\ summary transcript) is
given in Table~\ref{tab:qa_detail}.

\begin{table}[H]
\centering
\caption{Per-video qualitative Q\&A rating and observations (summary vs.\ original).}
\label{tab:qa_detail}
\scriptsize
\renewcommand{\arraystretch}{1.35}
\begin{tabular}{@{}p{1.5cm} p{1.8cm} p{10.5cm}@{}}
\toprule
\textbf{Video} & \textbf{Rating (summ)} & \textbf{Qualitative observation} \\
\midrule
video\_01 & \textbf{Incorrect}
  & Original: US--Israel--Iran news, political leaders.
    Summary: describes a \textbf{completely different topic}
    (e.g.\ antibiotics / comedy sketch) --- severe content-selection error. \\
video\_02 & \textbf{Accurate}
  & Both responses cover Israel/Iran, energy, US social issues;
    the summary version additionally captures an extremist-movement story
    in the US --- \textbf{topic matches}. \\
video\_03 & \textbf{Partial}
  & Original: football commentary with specific teams and players.
    Summary: correct sports direction (FA Cup, England--Brazil) but
    \textbf{missing} the specific team/player details present in the original. \\
video\_04 & \textbf{Accurate}
  & Both cover AI scenario planning and control risks ---
    \textbf{well aligned}. \\
video\_05 & \textbf{Accurate}
  & Original: neural electronics, pathology.
    Summary: BCIs, biomaterials --- \textbf{matches}, even more precise
    technically. \\
video\_06 & \textbf{Partial}
  & Original: Middle East, Iran/US/Russia, oil.
    Summary: correct geopolitical direction, but \textbf{adds} a
    ``Vietnamese traffic accident'' detail \textbf{absent} from the original
    (hallucination / context noise). \\
video\_07 & \textbf{Partial}
  & Original: Iran/Russia, Ukraine, energy, EU, Slovakia.
    Summary: correct on economy/China, but \textbf{adds}
    ``Vietnam economic competition'' which is \textbf{unrelated}. \\
video\_08 & \textbf{Partial}
  & Original: US--Iran, border, Khamenei.
    Summary: correct on Iran/Trump/China, but \textbf{adds} Chinese real
    estate and New York protests --- \textbf{off-topic} peripheral content. \\
video\_09 & \textbf{Partial}
  & Original: multiple domestic and international events.
    Summary: international part \textbf{approximately correct} but
    \textbf{misinterprets} some points; domestic content is \textbf{missing} ---
    consistent with only $\approx$9.5\% of transcript segments retained
    (Table~\ref{tab:vi_seg}). \\
video\_10 & \textbf{Accurate}
  & Both cover Middle East maritime issues, security, rescue operations ---
    \textbf{key topics preserved}. \\
\bottomrule
\end{tabular}
\end{table}

\noindent
\textbf{Interpretation.} Q\&A quality on the summary does \emph{not} track
CLIP search quality: when the transcript loses segment density or topic
fidelity, the LLM still produces fluent output that is \textbf{unfaithful}
to the original video. The two extremes are \textbf{video\_01}
(complete topic drift) and \textbf{video\_09} (severe information loss
due to only 9.5\% of segments retained).

\subsection{Limitations and Recommendations}

\begin{itemize}[leftmargin=1.8em]
  \item Sample size: 10 videos; one LLM; one question template---results do
        not generalize across all language and content domains.
  \item CLIP top-1 score does not substitute for human-labeled relevance
        judgment.
  \item \textbf{Recommendations:} (i) enforce a minimum segment count
        threshold before trusting Q\&A; (ii) check topic consistency between
        original and summary transcripts; (iii) verify audio coverage when an
        anomalous drop occurs.
\end{itemize}

% =============================================================================
\section{Conclusion}
% =============================================================================

This report has presented a \textbf{multimodal BiLSTM with Temporal Attention}
using \textbf{2432-dimensional} feature vectors, trained with supervised
learning and evaluated at summary ratio $r{=}0.35$. Ablation experiments
(Section~8.1) confirm the individual contribution of the audio-language branch,
bidirectional encoding, and attention mechanism, while the literature comparison
(Section~8.2) contextualizes the results within known benchmarks.

On the application side, \textbf{VideoIntel} demonstrates that \textbf{CLIP
semantic search} on video summaries suffers \textbf{negligible degradation} at
the defined metric level, while \textbf{question-answering quality} depends
critically on the \textbf{density and topical accuracy of the transcript}. This
reflects the nature of video summarization as \textbf{lossy compression} and
highlights the need for \textbf{downstream task validation} before real-world
deployment.

% =============================================================================
\section*{Appendix: Ablation Experiments}
% =============================================================================

All ablation variants (visual-only BiLSTM, full multimodal BiLSTM, no
attention, UniLSTM) have been reported with numerical results in
Section~\ref{tab:ablation_results}. Extensions to new domains or downstream
tasks (classification, captioning on summarized video) are left as future work
beyond the scope of this report.

% =============================================================================
%  References
% =============================================================================
\newpage
\bibliographystyle{ieeetr}

\begin{thebibliography}{99}

\bibitem{apostolidis2021pgl}
E.~Apostolidis, K.~Mekkas, A.~I.~Patras, and I.~Kompatsiaris,
``PGL-SUM: A Novel Pointer-Generator Network for Extractive Video
Summarization,'' in \textit{Proc. IEEE Int. Symp. Multimedia (ISM)}, 2021.

\bibitem{gygli2014}
M.~Gygli, H.~Grabner, H.~Riemenschneider, and L.~Van Gool,
``Creating summaries from user videos,''
in \textit{ECCV}, vol.~8695, 2014, pp.~505--520.

\bibitem{song2015}
Y.~Song, J.~Vallmitjana, A.~Stent, and A.~Jaimes,
``TVSum: Summarizing web videos using titles,''
in \textit{CVPR}, 2015, pp.~5179--5187.

\bibitem{fajtl2018}
J.~Fajtl, H.~S.~Sokeh, V.~Argyriou, D.~Monekosso, and P.~Remagnino,
``Summarizing Videos with Attention (VASNet),''
in \textit{ACCV}, 2018.

\bibitem{ghauri2021}
S.~A.~Ghauri, F.~S.~Khan, S.~W.~Zamir, J.~Hayat, and M.~Shah,
``MSVA: Multi-Stage Aggregated Transformer for Video Summarization,''
\textit{arXiv:2104.11530}, 2021.

\bibitem{mahasseni2017}
B.~Mahasseni, M.~Lam, and S.~Todorovic,
``Unsupervised Video Summarization with Adversarial LSTM Networks (SUM-GAN),''
in \textit{CVPR}, 2017.

\bibitem{zhang2016lstm}
K.~Zhang, W.-L.~Chao, F.~Sha, and K.~Grauman,
``Video Summarization with Long Short-Term Memory (vsLSTM, dppLSTM),''
in \textit{ECCV}, 2016.

\bibitem{zhang2019}
K.~Zhang, K.~Grauman, and F.~Sha,
``Retrospective Encoders for Video Summarization (Cycle-SUM),''
in \textit{AAAI}, 2019.

\bibitem{zhu2021dsnet}
W.~Zhu, J.~Lu, J.~Li, and J.~Zhou,
``DSNet: A Flexible Detect-to-Summarize Network for Video Summarization,''
\textit{IEEE Trans. Image Process.}, 2021.

\bibitem{zhu2022}
Y.~Zhu \textit{et al.},
``MHSCNet: Multi-Head Self-Attention Based Convolutional Network for Video
Summarization,'' \textit{arXiv:2204.08352}, 2022.

\bibitem{he2016}
K.~He, X.~Zhang, S.~Ren, and J.~Sun,
``Deep Residual Learning for Image Recognition,''
in \textit{CVPR}, 2016.

\bibitem{hochreiter1997}
S.~Hochreiter and J.~Schmidhuber,
``Long Short-Term Memory,''
\textit{Neural Computation}, vol.~9, no.~8, pp.~1735--1780, 1997.

\bibitem{radford2021clip}
A.~Radford \textit{et al.},
``Learning Transferable Visual Models From Natural Language Supervision
(CLIP),''
in \textit{ICML}, 2021.

\bibitem{reimers2019}
N.~Reimers and I.~Gurevych,
``Sentence-BERT: Sentence Embeddings using Siamese BERT-Networks,''
in \textit{EMNLP}, 2019.

\bibitem{radford2023whisper}
A.~Radford \textit{et al.},
``Robust Speech Recognition via Large-Scale Weak Supervision (Whisper),''
in \textit{ICML}, 2023.

\bibitem{ollama}
Ollama Developers,
\textit{Ollama: Run Large Language Models Locally},
\url{https://ollama.com}.

\bibitem{ffmpeg}
FFmpeg Developers,
\textit{FFmpeg: A Complete, Cross-Platform Solution to Record, Convert and
Stream Audio and Video},
\url{https://ffmpeg.org}.

\bibitem{haq2021}
H.~B.~Haq, M.~Asif, M.~B.~Ahmad, R.~Ashraf, and T.~Mahmood,
``Video summarization techniques: A review,''
\textit{Math. Probl. Eng.}, vol.~2021, pp.~1--17, 2021.

\bibitem{money2015}
A.~G.~Money and H.~Agius,
``A survey on video summarization techniques,''
\textit{Int. J. Comput. Appl.}, vol.~118, no.~11, pp.~25--31, 2015.

\bibitem{psallidas2021}
T.~Psallidas, P.~Koromilas, T.~Giannakopoulos, and E.~Spyrou,
``Multimodal summarization of user-generated videos,''
\textit{Appl. Sci.}, vol.~11, no.~11, p.~5260, 2021.

\bibitem{hu2023}
S.~Hu, Z.~Liu, J.~Liu, and Z.~Guo,
``Video summarization based on feature fusion and data augmentation,''
\textit{Computers}, vol.~12, no.~9, p.~186, 2023.

\bibitem{kadam2022}
P.~Kadam \textit{et al.},
``Recent challenges and opportunities in video summarization with machine
learning algorithms,''
\textit{IEEE Access}, vol.~10, pp.~122762--122785, 2022.

\bibitem{apostolidis2021survey}
E.~Apostolidis, E.~Adamantidou, A.~I.~Metsai, V.~Mezaris, and I.~Patras,
``Video summarization using deep neural networks: A survey,''
\textit{Proc. IEEE}, vol.~109, no.~11, pp.~1838--1863, 2021.

\bibitem{bertasius2021}
G.~Bertasius, H.~Wang, and L.~Torresani,
``Is Space-Time Attention All You Need for Video Understanding? (TimeSformer),''
in \textit{ICML}, 2021.

\bibitem{liu2022swin}
Z.~Liu \textit{et al.},
``Video Swin Transformer,''
in \textit{CVPR}, 2022.

\bibitem{carreira2017}
J.~Carreira and A.~Zisserman,
``Quo Vadis, Action Recognition? A New Model and the Kinetics Dataset (I3D),''
in \textit{CVPR}, 2017.

\bibitem{kingma2015}
D.~P.~Kingma and J.~Ba,
``Adam: A Method for Stochastic Optimization,''
in \textit{ICLR}, 2015.

\bibitem{loshchilov2019}
I.~Loshchilov and F.~Hutter,
``Decoupled Weight Decay Regularization (AdamW),''
in \textit{ICLR}, 2019.

\end{thebibliography}

\end{document}